\scanpage{107}
Soit $J$ un ensemble ordonné, on s'intéresse aux plongements d'ensembles ordonnés
\[
(150)\qquad J\mathrel{\overset{\varphi}{\hookrightarrow}}\Delta_{\unsure{x|\alpha}}\unsure{\xi}\simeq\mathcal P(I(\unsure{x|\alpha}))^\circ
\]
Un tel plongement donne une ``fonction codimension'' sur $J$, en composant, \inserted{\unsure{soit} $d$,} et \deleted{\unsure{dire} que $\varphi$ est un plongement d'\unsure{ens.} \unsure{ordonnés} signifie que}\inserted{on suppose donné déjà $d$ d'avance sur $J$} satisfaisant
\[
(151)\qquad\text{Si }x<y\quad\text{on a}\quad\unsure{d|c}(x)<c(y),
\]
et on se borne aux $\varphi$ qui sont des plongts ordonnés, compatibles aux foncts codim sur $J$ et sur $\Delta_{\unsure{x|\alpha}}$.
\deleted{On suppose de plus que $J$ \unsure{a} un plus petit élément (\unsure{ce} qui correspond donc à une codim. \unsure{maximale|minimale} \deleted{\unsure{qui est}} soit $d(J)$) — quitte à le rajouter, s'il n'y est déjà!}

Considérons, \inserted{pour tout $\unsure{q|d|p}\in\mathbb N^*$,} la multiplicité analytique formée des couples $(x,V)$, où $x$ est un pt de $D$ et $V$ un germe de sous-multi-
plicité de $X$ attaché à $D$, de codim $\unsure{q|d|p}$. On \deleted{\unsure{voit}} voit aisément que ce sont \inserted{bien} là les points d'une multiplicité analytique que je note $D_{\unsure{q|d|p}}$, et on a un morphisme
\margin{106}
